\chapter{The Kind One}

The poem stayed with them long after it was spoken.

Neither Father Handy nor Tibor knew who had written it. The words were older than their dictionaries, older than certainty. They had agreed on its meaning by instinct rather than knowledge, and even then only partially. It spoke of light and green places, of going somewhere together—yet whether that journey was toward life or away from it remained unclear.

That uncertainty frightened them.

And soothed them.

It was the kind of truth that could not be resolved, only returned to. A balm precisely because it resisted explanation. Ely, watching from the corner, sneered. She disliked anything that could not be pinned down and labeled.

Father Handy needed power now. Mekkis, he thought. Not the kind that came from engines or batteries, but the kind that pressed downward from above, invisible and absolute. Even the Servants of Wrath agreed with the old Christians on that much: whatever ruled this world came from beyond the stars.

Yet he distrusted himself for clinging to fragments—half-understood poems, obsolete maps, scraps of meaning from a dead civilization. Perhaps degeneration wasn’t just historical. Perhaps it had settled inside them.

The Dominus McComas arrived that afternoon, heavy and coarse, a man who looked as though he consumed the world rather than lived in it. His teeth were wrong—too aggressive, too visible, as if they were tools rather than parts of a body.

“Carleton Lufteufel was a bastard,” McComas said cheerfully. “As a man, I mean.”

Father Handy said nothing. He spread old gas-station maps across the table, the paper soft with age. Roads that no longer existed still promised connection.

McComas laughed, spat, scratched his gums with a matchstick. “So your artist is going all the way west? On that cow-cart? He won’t make it. That’s mercy, really.”

Father Handy ignored him and thought instead of the war.

Not the missiles. Not the planes. But the gob—the weapon that didn’t strike a place, only the sky. It poisoned the atmosphere itself. No defense could stop it. When Lufteufel triggered it from orbit, the world didn’t explode. It sickened.

People thinned. Blood failed. Bodies unraveled slowly.

Other weapons had been clever, even elegant—defensive aircraft that hovered for months, pilots sealed into self-sustaining suits, machines sacrificing three lives to save millions below. But cleverness had limits. Everything did.

What haunted Father Handy wasn’t the scale of death, but the intimacy of certain weapons. He thought of the nerve gas. The way organs consumed themselves. Evil didn’t have to be efficient to be complete.

McComas mocked the mural. “Paint him fat and ugly,” he said. “That’s all people need.”

But Father Handy had seen the photograph. The god had looked human. That was the problem.

Then Lurine Rae arrived.

She entered lightly, as if gravity negotiated with her. Red hair, freckles, small bones. A woman who seemed capable of flight. She unsettled Father Handy in a way doctrine never had.

She disliked McComas openly. He disliked her instinctively.

“I’m joining the Christian church,” she said calmly, over coffee.

McComas roared with laughter. “What church?”

“The kind one,” she said.

Father Handy felt the ground shift.

Christianity was supposed to be finished. Its god had been exposed as a lie—or worse, a mask. To return to it now felt like denial. Escape.

They argued. About sin. About knowledge. About whether death was release or defeat.

Lurine refused to worship a former government official as a god. She said it plainly. That mattered.

“But he lives,” Father Handy said.

That silenced her.

Later, when McComas had left, Father Handy asked her why.

“I like kind people,” she said simply.

That answer terrified him more than any theology.

He tried to explain. Tried to warn her that belief in goodness was a flight from reality. That the world had already proven what ruled it.

She listened. She did not argue. But she did not yield.

“He was a man,” she said finally. “Not a god.”

Father Handy had no reply that did not circle back to fear.

Outside, the world remained broken. Roads faded into dirt. Maps lied. People survived by growing food, not preserving data. Knowledge had not saved them.

And yet Lurine sat there, calm and determined, believing in something gentler.

It made no sense.

And that, more than anything else, unsettled him.

